\section{Main Theorem}

\begin{theorem}[Conditional static mean-response-rank theorem]
\label{thm:main}
Under Assumptions~\ref{assump:parameter-regime}--\ref{assump:mean-response-rank},
there is a basis \(\psi_1,\ldots,\psi_{r_A}\) of \(V_A\), chosen once from
the fixed domain, learner, seed law, horizon, tolerance, complete-response
interface, and a fixed basis choice, such that the deterministic map
\[
\varphi_A(x):=(\psi_1(x),\ldots,\psi_{r_A}(x))
\in\mathbb R^{r_A}
\]
is independent of every \(\mathcal D\), target, tolerance-valid response
policy, and realized learner seed.  For every \(h\in\mathcal H\), there is a
weight \(w_h\in\mathbb R^{r_A}\) satisfying
\[
h(x)\langle w_h,\varphi_A(x)\rangle
\ge 1-2\varepsilon=\rho>\frac12>0
\qquad\text{for every }x\in\mathcal X.
\]
Consequently,
\[
\operatorname{dc}(\mathcal H)
\le r_A
\le B\bigl(1+m/\tau^2\bigr)^k.
\]
The conclusion is deterministic, pointwise, and fixed-horizon, with no hidden
constants or domain-cardinality dependence.  If \(\mathcal X=\varnothing\)
or \(\mathcal H=\varnothing\), then \(\operatorname{dc}(\mathcal H)=0\).
If both are nonempty, then \(r_A\ge1\).  At \(m=0\), the bound specializes
to \(r_A\le B\); at \(\varepsilon=0\), the signed margin is exactly \(1\);
and \(B=k=1\) gives \(r_A\le1+m/\tau^2\).
\end{theorem}

The theorem is conditional: Assumption~\ref{assump:mean-response-rank} is an
additional static protocol-family certificate.  The argument does not derive
that certificate from \(m\) and \(\tau\) alone and does not establish the
unconditional universal linear dimension bound.  Those source-level questions
remain open.
